\documentclass[12pt,a4paper]{article}

% ---------- Paquetes ----------
\usepackage[spanish]{babel}
\usepackage[T1]{fontenc}
\usepackage[utf8]{inputenc}
\usepackage{lmodern}
\usepackage[a4paper,margin=2.5cm]{geometry}
\usepackage{amsmath,amssymb,amsthm,mathtools}
\usepackage{physics}
\usepackage{braket} % Adicional para Dirac
\usepackage{bm}
\usepackage{siunitx}
\usepackage{booktabs}
\usepackage{array}
\usepackage{enumitem}
\usepackage{graphicx}
\usepackage{subcaption}
\usepackage{float}
\usepackage{caption}
\usepackage{hyperref}
\usepackage{cleveref}
\usepackage{xcolor}
\usepackage{listings}

% ---------- Configuración ----------
\hypersetup{
  colorlinks=true,
  linkcolor=blue!60!black,
  urlcolor=blue!60!black,
  citecolor=blue!60!black
}

\setlist[itemize]{topsep=4pt,itemsep=2pt}
\setlist[enumerate]{topsep=4pt,itemsep=2pt}

\captionsetup{font=small,labelfont=bf}

\lstdefinestyle{pythonstyle}{
  language=Python,
  basicstyle=\ttfamily\small,
  keywordstyle=\color{blue!70!black}\bfseries,
  commentstyle=\color{green!40!black},
  stringstyle=\color{orange!70!black},
  numbers=left,
  numberstyle=\tiny\color{gray},
  stepnumber=1,
  numbersep=8pt,
  showstringspaces=false,
  breaklines=true,
  frame=single,
  rulecolor=\color{black!20}
}

\lstset{style=pythonstyle}

\newcommand{\vect}[1]{\bm{#1}}
\newcommand{\R}{\mathbb{R}}
\newcommand{\informeimg}[2]{%
  \IfFileExists{#1}{\includegraphics[width=#2]{#1}}{%
    \fbox{\parbox[c][5cm][c]{0.9\linewidth}{\centering
    Imagen pendiente:\\ \texttt{\detokenize{#1}}}}%
  }%
}

% ---------- Datos del informe ----------
\title{\textbf{Laboratorio 7}\\
\large Computación Cuántica}
\author{Edmil Jampier Saire Bustamante\\
\texttt{174449@unsaac.edu.pe}}
\date{2026-1 --- IF460AIN-COMPUTACION CUANTICA --- IN-203}

\begin{document}

\maketitle
\thispagestyle{empty}

\begin{center}
\begin{tabular}{ll}
\textbf{Curso:} & Computación Cuántica \\
\textbf{Docente:} & Raul Huillca Huallparimachi \\
\textbf{Comisión:} & IN-203
\end{tabular}
\end{center}

\vspace{1em}

\begin{abstract}
En este informe de laboratorio se desarrollan ejercicios teóricos sobre fundamentos matemáticos para Computación Cuántica. Se estudian conceptos clave como vectores de probabilidad, matrices estocásticas, notación de Dirac, validación de vectores de estado cuántico, ordenamiento de productos cartesianos y cálculo de productos tensoriales. Además de la resolución analítica, se incluye la implementación numérica en Python utilizando la librería NumPy, verificando de manera empírica cada una de las soluciones propuestas y visualizándolas gráficamente.
\end{abstract}

\tableofcontents
\newpage

\section{1. Reconocer vectores de probabilidad}

Sea $\Sigma = \{0,1,2,\ldots,9\}$ el conjunto de estados clásicos de un sistema que representa un solo dígito.

\textbf{¿Cuáles de los siguientes vectores son vectores de probabilidad?} Justifique su respuesta.

\begin{enumerate}
    \item $\frac{2}{5}|8\rangle+\frac{1}{5}\left(|3\rangle+|4\rangle+|5\rangle\right)$
    \item $\frac{3}{2}|0\rangle+\frac{1}{2}|1\rangle-|0\rangle$
    \item $\frac{1}{10}\sum_{j=1}^{9}|j\rangle$
    \item $(1-\sqrt{2})|7\rangle+\sqrt{2}|2\rangle$
    \item $\frac{1}{3}|3\rangle+\frac{1}{4}|4\rangle+\frac{5}{12}|5\rangle$
\end{enumerate}

\begin{figure}[H]
    \centering
    \informeimg{imagenes/ej1_prob_vectors.png}{0.75\textwidth}
    \caption{Visualización de vectores de probabilidad válidos.}
\end{figure}

\vspace{0.5em}
\textbf{Solución:}
Un vector es de probabilidad si la suma de sus componentes es $1$ y todas sus componentes son $\geq 0$.
\begin{itemize}
    \item \textbf{1. Sí.} $\frac{2}{5} + \frac{1}{5}(1+1+1) = \frac{2}{5} + \frac{3}{5} = 1$. Todas son $\geq 0$.
    \item \textbf{2. Sí.} El vector evaluado es $\left(\frac{3}{2}-1\right)|0\rangle + \frac{1}{2}|1\rangle = \frac{1}{2}|0\rangle + \frac{1}{2}|1\rangle$. La suma es $\frac{1}{2}+\frac{1}{2}=1$.
    \item \textbf{3. No.} La suma es $9 \times \frac{1}{10} = \frac{9}{10} \neq 1$.
    \item \textbf{4. No.} Contiene un componente negativo, ya que $1-\sqrt{2} < 0$.
    \item \textbf{5. Sí.} La suma es $\frac{1}{3} + \frac{1}{4} + \frac{5}{12} = \frac{4}{12} + \frac{3}{12} + \frac{5}{12} = \frac{12}{12} = 1$.
\end{itemize}

\section{2. Reconocer matrices estocásticas}

Las siguientes matrices son todas matrices de $3\times 3$.

\textbf{Seleccione todas aquellas que sean matrices estocásticas.}

\begin{enumerate}
    \item $|0\rangle\langle0|+|1\rangle\langle1|-|2\rangle\langle2|$
    \item $\begin{pmatrix} \frac{1}{3} & \frac{1}{3} & \frac{1}{3}\\ \frac{1}{4} & \frac{5}{12} & \frac{1}{3}\\ \frac{5}{12} & \frac{1}{3} & \frac{1}{4} \end{pmatrix}$
    \item $\begin{pmatrix} 0 & 1 & 0\\ 0 & 0 & 1\\ 1 & 0 & 0 \end{pmatrix}$
    \item $\frac{1}{3}\left(|0\rangle+|1\rangle+|2\rangle\right)\left(\langle0|+\langle1|+\langle2|\right)$
    \item $\begin{pmatrix} 1 & 1 & 0\\ 0 & 0 & 1\\ 1 & 0 & 0 \end{pmatrix}$
    \item $\begin{pmatrix} 1 & 0 & 0\\ 1 & 0 & 0\\ 0 & 1 & 0 \end{pmatrix}$
\end{enumerate}

\begin{figure}[H]
    \centering
    \informeimg{imagenes/ej2_stochastic_matrices.png}{0.9\textwidth}
    \caption{Matrices estocásticas identificadas representadas como mapas de calor.}
\end{figure}

\vspace{0.5em}
\textbf{Solución:}
Para ser una matriz estocástica (por filas), todas las entradas deben ser $\geq 0$ y la suma de cada fila debe ser $1$.
\begin{itemize}
    \item \textbf{1. No.} Contiene una entrada negativa ($-1$).
    \item \textbf{2. Sí.} Todas las entradas son $\geq 0$. Suma de Fila 1: $\frac{1}{3}+\frac{1}{3}+\frac{1}{3}=1$. Fila 2: $\frac{1}{4}+\frac{5}{12}+\frac{1}{3}=\frac{3+5+4}{12}=1$. Fila 3: $\frac{5}{12}+\frac{1}{3}+\frac{1}{4}=1$.
    \item \textbf{3. Sí.} Es una matriz de permutación, la suma por filas es $1$.
    \item \textbf{4. Sí.} Todas las entradas son $\frac{1}{3}$, la suma por filas es $3 \times \frac{1}{3} = 1$.
    \item \textbf{5. No.} La suma de la primera fila es $1+1+0=2 \neq 1$.
    \item \textbf{6. Sí.} Todas las sumas por fila son $1$.
\end{itemize}

\section{3. Convertir matriz a notación de Dirac}

Sean $|0\rangle$ y $|1\rangle$ los vectores base estándar para un solo qubit:
$$|0\rangle= \begin{pmatrix} 1\\ 0 \end{pmatrix}, \qquad |1\rangle= \begin{pmatrix} 0\\ 1 \end{pmatrix}$$

¿Cuál de las siguientes es una representación correcta de la matriz
$$M= \begin{pmatrix} 2 & -1\\ -9 & 3 \end{pmatrix}$$
en notación de Dirac? \textbf{Demuestre todas las opciones.}

\begin{enumerate}
    \item $M= 2(2|0\rangle-|1\rangle) (-9\langle0|+3\langle1|)$
    \item $M= 2|0\rangle\langle0| -|0\rangle\langle1| -9|1\rangle\langle0| +3|1\rangle\langle1|$
    \item $M= 2|1\rangle\langle1| -|0\rangle\langle1| -9|1\rangle\langle0| +3|0\rangle\langle0|$
    \item $M= 2|0\rangle\langle0| -|1\rangle\langle0| -9|0\rangle\langle1| +3|1\rangle\langle1|$
    \item $M= 2|00\rangle\langle00| -|01\rangle\langle01| -9|10\rangle\langle10| +3|11\rangle\langle11|$
\end{enumerate}

\vspace{0.5em}
\textbf{Solución:}
La matriz $M$ se puede descomponer en notación de Dirac sumando cada componente $M_{ij} |i\rangle\langle j|$:
$$M = M_{00}|0\rangle\langle0| + M_{01}|0\rangle\langle1| + M_{10}|1\rangle\langle0| + M_{11}|1\rangle\langle1|$$
Reemplazando con los valores de la matriz:
$$M = 2|0\rangle\langle0| - 1|0\rangle\langle1| - 9|1\rangle\langle0| + 3|1\rangle\langle1|$$
Por lo tanto, la opción correcta es la \textbf{2}. Las opciones 1, 3, 4 y 5 al ser desarrolladas no coinciden con esta expresión matricial.

\section{4. Reconocer vectores de estados cuánticos}

Sea $\Sigma=\{\#,b,c\}$ el conjunto de estados clásicos de un sistema.

\textbf{¿Cuáles de los siguientes son vectores de estados cuánticos?} Seleccione todas las que correspondan.

\begin{enumerate}
    \item $\frac{1}{2} \left( \sqrt{6}|\#\rangle -\sqrt{2}|b\rangle +2i|c\rangle \right)$
    \item $\sqrt{\frac{i}{3}}|c\rangle +\sqrt{\frac{1}{2}}|\#\rangle +\sqrt{\frac{1}{6}}|b\rangle$
    \item $\frac{1}{2} \left( \sqrt{2}|c\rangle -\sqrt{\frac{2}{3}}|\#\rangle +\sqrt{\frac{2i}{3}}|b\rangle \right)$
    \item $\frac{i-\sqrt{3}}{2\sqrt{3}}|c\rangle +\sqrt{\frac{1-i}{12}}|b\rangle +\sqrt{\frac{1}{2}}|\#\rangle$
\end{enumerate}

\begin{figure}[H]
    \centering
    \informeimg{imagenes/ej4_quantum_states.png}{0.75\textwidth}
    \caption{Distribución de probabilidad de los estados cuánticos válidos.}
\end{figure}

\vspace{0.5em}
\textbf{Solución:}
Para que un vector sea de estado cuántico, su norma al cuadrado debe ser $1$ ($\sum |\alpha_i|^2 = 1$).
\begin{itemize}
    \item \textbf{1. No.} Norma$^2 = \frac{1}{4}\left(|\sqrt{6}|^2 + |-\sqrt{2}|^2 + |2i|^2\right) = \frac{1}{4}(6 + 2 + 4) = 3 \neq 1$.
    \item \textbf{2. Sí.} Norma$^2 = \left|\sqrt{\frac{i}{3}}\right|^2 + \left|\sqrt{\frac{1}{2}}\right|^2 + \left|\sqrt{\frac{1}{6}}\right|^2 = \frac{1}{3} + \frac{1}{2} + \frac{1}{6} = \frac{2+3+1}{6} = 1$.
    \item \textbf{3. No.} Norma$^2 = \frac{1}{4}\left(|\sqrt{2}|^2 + \left|-\sqrt{\frac{2}{3}}\right|^2 + \left|\sqrt{\frac{2i}{3}}\right|^2\right) = \frac{1}{4}\left(2 + \frac{2}{3} + \frac{2}{3}\right) = \frac{10}{12} = \frac{5}{6} \neq 1$.
    \item \textbf{4. No.} Norma$^2 = \left|\frac{i-\sqrt{3}}{2\sqrt{3}}\right|^2 + \left|\sqrt{\frac{1-i}{12}}\right|^2 + \left|\sqrt{\frac{1}{2}}\right|^2 = \frac{1+3}{12} + \frac{\sqrt{2}}{12} + \frac{1}{2} = \frac{5}{6} + \frac{\sqrt{2}}{12} \neq 1$.
\end{itemize}

\section{5. Orden alfabético de productos cartesianos}

Sea $\Gamma=\{1,2,3,4\}$ ordenado como está escrito, con $1$ primero.

\textbf{¿Cuál es el elemento número 12 del conjunto $\Gamma\times\Gamma$ ordenado alfabéticamente?}

\vspace{0.5em}
\textbf{Solución:}
El conjunto $\Gamma \times \Gamma$ tiene $4 \times 4 = 16$ elementos ordenados de la siguiente forma:
$$(1,1), (1,2), (1,3), (1,4),$$
$$(2,1), (2,2), (2,3), (2,4),$$
$$(3,1), (3,2), (3,3), (3,4), \dots$$
Como cada número inicial agrupa $4$ elementos, el elemento número 12 es el último del tercer grupo, es decir, \textbf{$(3,4)$}.

\section{6. Calcular productos tensoriales}

Se tienen los siguientes vectores:
$$|u\rangle= \begin{pmatrix} -2\\ 1 \end{pmatrix}, \qquad |v\rangle= \begin{pmatrix} 1\\ 0\\ 1 \end{pmatrix}$$
$$|w\rangle= \begin{pmatrix} 4\\ 1 \end{pmatrix}, \qquad |x\rangle= \begin{pmatrix} \frac{1}{2}\\ \frac{1}{2}\\ \frac{1}{2} \end{pmatrix}$$

Calcular el siguiente producto tensorial:
$$|u\rangle\otimes|v\rangle -\frac{2}{3} |w\rangle\otimes|x\rangle$$

\vspace{0.5em}
\textbf{Solución:}
Calculamos primero el producto tensorial $|u\rangle\otimes|v\rangle$:
$$|u\rangle\otimes|v\rangle = \begin{pmatrix} -2 \\ 1 \end{pmatrix} \otimes \begin{pmatrix} 1 \\ 0 \\ 1 \end{pmatrix} = \begin{pmatrix} -2(1) \\ -2(0) \\ -2(1) \\ 1(1) \\ 1(0) \\ 1(1) \end{pmatrix} = \begin{pmatrix} -2 \\ 0 \\ -2 \\ 1 \\ 0 \\ 1 \end{pmatrix}$$

Luego calculamos $\frac{2}{3}|w\rangle\otimes|x\rangle$:
$$\frac{2}{3}|w\rangle\otimes|x\rangle = \frac{2}{3} \begin{pmatrix} 4 \\ 1 \end{pmatrix} \otimes \begin{pmatrix} 1/2 \\ 1/2 \\ 1/2 \end{pmatrix} = \frac{2}{3} \begin{pmatrix} 2 \\ 2 \\ 2 \\ 1/2 \\ 1/2 \\ 1/2 \end{pmatrix} = \begin{pmatrix} 4/3 \\ 4/3 \\ 4/3 \\ 1/3 \\ 1/3 \\ 1/3 \end{pmatrix}$$

Restando ambos resultados obtenemos:
$$|u\rangle\otimes|v\rangle - \frac{2}{3} |w\rangle\otimes|x\rangle = \begin{pmatrix} -2 - 4/3 \\ 0 - 4/3 \\ -2 - 4/3 \\ 1 - 1/3 \\ 0 - 1/3 \\ 1 - 1/3 \end{pmatrix} = \begin{pmatrix} -10/3 \\ -4/3 \\ -10/3 \\ 2/3 \\ -1/3 \\ 2/3 \end{pmatrix}$$

\end{document}
